\text{UNIT 9: TRANSFORMATIONS}

\text{=================================}

\text{TOPIC 1: TRANSLATIONS}

\text{Rule form: (x,y) \to (x+a, y+b)}

\text{Vector form: \langle a,b \rangle means}
\text{right a, up b (negatives flip)}

\text{Add a to x, add b to y}

\text{Example Q1}

\text{D(6,1) with (x,y) \to (x-11, y+4)}

\text{D'=(6-11, 1+4)=(-5,5)}

\text{E(8,0) \to E'=(-3,4)}

\text{F(5,-6) \to F'=(-6,-2)}

\text{G(3,-5) \to G'=(-8,-1)}

\text{Example Q3 vector \langle -4,0 \rangle}

\text{P(2,7) \to P'=(-2,7)}

\text{Q(8,-3) \to Q'=(4,-3)}

\text{R(-3,-2) \to R'=(-7,-2)}

\text{Finding the rule from image:}

\text{Subtract pre from image}

\text{a = x'-x, b = y'-y}


\text{=================================}

\text{TOPIC 2: REFLECTIONS}

\text{Over x-axis: (x,y) \to (x,-y)}

\text{Over y-axis: (x,y) \to (-x,y)}

\text{Over y=x: (x,y) \to (y,x)}

\text{Over y=-x: (x,y) \to (-y,-x)}

\text{Over y=k (horiz line):}
\text{(x,y) \to (x, 2k-y)}

\text{Over x=h (vert line):}
\text{(x,y) \to (2h-x, y)}

\text{Example: (3,5) over x-axis}
\text{\to (3,-5)}

\text{Example: (-7,3) over x-axis}
\text{\to (-7,-3)}

\text{Reverse: if image (-7,3),}
\text{preimage N=(-7,-3)}


\text{=================================}

\text{TOPIC 3: ROTATIONS ABOUT ORIGIN}

\text{Counterclockwise (CCW):}

\text{90 CCW: (x,y) \to (-y,x)}

\text{180: (x,y) \to (-x,-y)}

\text{270 CCW: (x,y) \to (y,-x)}

\text{Clockwise (CW):}

\text{90 CW = 270 CCW: (y,-x)}

\text{270 CW = 90 CCW: (-y,x)}

\text{Example Q16: 90 CCW}

\text{C(2,4) \to C'=(-4,2)}

\text{D(6,7) \to D'=(-7,6)}

\text{E(7,3) \to E'=(-3,7)}

\text{Example Q19: 90 CW}

\text{P(1,3) \to P'=(3,-1)}

\text{Q(4,8) \to Q'=(8,-4)}

\text{R(8,1) \to R'=(1,-8)}

\text{ROTATION ABOUT POINT (a,b):}

\text{Step 1: subtract center}
\text{(x-a, y-b)}

\text{Step 2: apply rotation rule}

\text{Step 3: add center back}

\text{Example Q20: 90 CCW about (2,1)}

\text{D(-7,8): shift \to (-9,7)}
\text{rotate \to (-7,-9)}
\text{shift back \to (-5,-8)}

\text{So D'=(-5,-8)}


\text{=================================}

\text{TOPIC 4: DILATIONS}

\text{Centered at origin:}

\text{(x,y) \to (kx, ky)}

\text{k>1 enlargement, 0<k<1 reduction}

\text{Example Q22: k=2 at origin}

\text{A(-7,1) \to A'=(-14,2)}

\text{B(-5,6) \to B'=(-10,12)}

\text{C(-1,2) \to C'=(-2,4)}

\text{Example Q23: k=3/4 at origin}

\text{L(-12,0) \to L'=(-9,0)}

\text{M(-4,-4) \to M'=(-3,-3)}

\text{N(-4,-8) \to N'=(-3,-6)}

\text{O(-12,-12) \to O'=(-9,-9)}

\text{DILATION ABOUT POINT (a,b):}

\text{(x,y) \to (k(x-a)+a, k(y-b)+b)}

\text{Example Q24: k=1/3 at (5,2)}

\text{J(8,11):}
\text{x'=(1/3)(8-5)+5=1+5=6}
\text{y'=(1/3)(11-2)+2=3+2=5}
\text{J'=(6,5)}

\text{K(14,11) \to K'=(8,5)}

\text{L(14,-4) \to L'=(8,0)}

\text{M(8,-4) \to M'=(6,0)}

\text{FIND SCALE FACTOR k:}

\text{k = image length / preimage length}

\text{or k = (image coord - center) /}
\text{(pre coord - center)}

\text{FIND CENTER:}

\text{Draw lines through corresponding}
\text{points; they meet at center}


\text{=================================}

\text{TOPIC 5: MIXED (REVERSE)}

\text{Work BACKWARDS from image}

\text{Q30: P'(-2,-5), rule (x+1,y+2)}
\text{P=(-2-1, -5-2)=(-3,-7)}

\text{Q31: A'(4,-1), vector \langle 0,-6 \rangle}
\text{A=(4-0, -1+6)=(4,5)}

\text{Q32: N'(-7,3) reflect x-axis}
\text{N=(-7,-3)}

\text{Q33: R'(2,-9) reflect y=x}
\text{R=(-9,2) (swap back)}

\text{Q34: D'(3,5) after 270 CCW}
\text{Inverse of 270 CCW is 90 CCW}
\text{So D = (-y',x') pattern reversed}
\text{270 CCW: (x,y)\to(y,-x)}
\text{Set (y,-x)=(3,5): y=3, x=-5}
\text{D=(-5,3)}

\text{Q35: V'(-6,-2), k=1/2 at origin}
\text{V=(-6 \div 1/2, -2 \div 1/2)}
\text{V=(-12,-4)}


\text{=================================}

\text{TOPIC 6: SEQUENCES}

\text{Apply transformations IN ORDER}

\text{Example Q36: D(1,0), E(7,2),}
\text{F(8,-1), G(2,-3)}

\text{Step a: (x-8, y-3)}
\text{D\to(-7,-3), E\to(-1,-1)}
\text{F\to(0,-4), G\to(-6,-6)}

\text{Step b: reflect x-axis (y\to-y)}
\text{D'=(-7,3), E'=(-1,1)}
\text{F'=(0,4), G'=(-6,6)}

\text{Example Q37: M(-2,0), N(0,-1),}
\text{P(-3,-2)}

\text{Step a: 90 CCW (x,y)\to(-y,x)}
\text{M\to(0,-2), N\to(1,0), P\to(2,-3)}

\text{Step b: dilate k=3 at origin}
\text{M'=(0,-6), N'=(3,0), P'=(6,-9)}


\text{=================================}

\text{TOPIC 7: SYMMETRY}

\text{LINE symmetry: fold to match}

\text{POINT symmetry: 180 rotation}
\text{looks the same}

\text{ROTATIONAL symmetry: rotates}
\text{less than 360 to match}

\text{Letter P: none}

\text{Letter plus/X shapes: all three}

\text{No-entry sign (circle + bar):}
\text{line, point, rotational}


\text{=================================}

\text{UNIT 10: CIRCLES}

\text{=================================}

\text{TOPIC 1: PARTS OF CIRCLE}

\text{Center: point in middle}

\text{Radius: center to edge}

\text{Diameter: edge to edge thru center}
\text{d = 2r}

\text{Chord: segment with both ends}
\text{on the circle}

\text{Secant: line that cuts circle}
\text{(chord extended)}

\text{Tangent: line touching at 1 point}

\text{Central angle: vertex AT center}

\text{Inscribed angle: vertex ON circle}

\text{Minor arc: less than 180}

\text{Major arc: more than 180}
\text{(named with 3 letters)}

\text{Semicircle: exactly 180}


\text{=================================}

\text{TOPIC 2: AREA & CIRCUMFERENCE}

\text{C = 2\pi r = \pi d}

\text{A = \pi r^2}

\text{Example Q2: diameter 16.6 m}
\text{r = 8.3}
\text{C = 2\pi(8.3) = 52.15 m}
\text{A = \pi(8.3)^2 = 216.42 m^2}

\text{Example Q3: C=106.81 cm, find r}
\text{r = C / (2\pi) = 17.0 cm}

\text{Example Q4: A=95.03, find d}
\text{r = \sqrt{A/\pi} = 5.5}
\text{d = 11 ft}

\text{Example Q5: A=254.47, find C}
\text{r = \sqrt{A/\pi} = 9}
\text{C = 2\pi(9) = 56.55 in}

\text{Example Q6: C=30\pi, find A}
\text{r = 15}
\text{A = \pi(15)^2 = 225\pi m^2}


\text{=================================}

\text{TOPIC 3: CENTRAL ANGLES & ARCS}

\text{Central angle = arc measure}

\text{All arcs around circle = 360}

\text{Example Q7: given 79 and 53}
\text{(and figure with 5 arcs labeled)}

\text{Set up: sum all arcs = 360}

\text{Solve for unknown arcs}

\text{SOLVE FOR x with central angles:}

\text{Sum of all central angles = 360}

\text{Example Q8: 45 + (6x+33) + 28}
\text{+ other arcs = 360}

\text{Combine and solve}

\text{VERTICAL central angles are equal}


\text{=================================}

\text{TOPIC 4: ARC LENGTHS}

\text{Arc length = (arc/360) \cdot 2\pi r}

\text{or in radians: s = r\theta}

\text{Example Q10: r=15, arc=26}
\text{VW length = (26/360) \cdot 2\pi(15)}
\text{= (26/360)(30\pi)}
\text{= 6.81 cm}

\text{Example Q11: UXV}
\text{Need arc measure first}
\text{Then plug into formula}


\text{=================================}

\text{TOPIC 5: CHORDS & ARCS}

\text{CONGRUENT CHORDS \to CONGRUENT ARCS}

\text{If chords equal, arcs equal}

\text{Example Q12: DC=(12x+7),}
\text{CB=(18x-23), chords congruent}

\text{12x+7 = 18x-23}
\text{30 = 6x}
\text{x = 5}

\text{arc DC = 12(5)+7 = 67}
\text{arc DAB = 360 - DC - CB}
\text{= 360 - 67 - 67 = 226}

\text{PERPENDICULAR FROM CENTER:}

\text{Radius perpendicular to chord}
\text{BISECTS the chord and its arc}

\text{Q13: GP=PH means chord bisected}

\text{If GA=17, full chord = 34}
\text{DB = 34}

\text{If two chords equidistant from}
\text{center, they are congruent}


\text{=================================}

\text{TOPIC 6: INSCRIBED ANGLES}

\text{INSCRIBED ANGLE THEOREM:}

\text{inscribed angle = (1/2) arc}

\text{arc = 2 \cdot inscribed angle}

\text{Same arc \to same inscribed angle}

\text{Angle in semicircle = 90}

\text{OPPOSITE angles in inscribed}
\text{quadrilateral sum to 180}

\text{Example Q14: arc JM=134}
\text{and angle MJN=59 (so arc MN=118)}

\text{a) mML = ?}
\text{b) JLK = (1/2) arc JK}
\text{c) JLM = (1/2) arc JM}
\text{= (1/2)(134) = 67}
\text{d) arc MJ = 2 \cdot angle MLJ}

\text{Example Q15: S=94, T=(9x-25)}
\text{inscribed quad: S+T=180}
\text{94 + 9x-25 = 180}
\text{9x = 111, x = 12.33}

\text{Q16: BCD=(7x+10), BAD=(19x-50)}
\text{opposite angles in quad: sum=180}
\text{7x+10+19x-50=180}
\text{26x = 220, x = 8.46}
\text{BAD = 19(8.46)-50 = 110.7}

\text{Q17: PQS=(6x+1), PRS=(34-5x)}
\text{same arc \to angles equal}
\text{6x+1 = 34-5x}
\text{11x = 33, x = 3}
\text{PS arc = 2(6(3)+1) = 38}

\text{Q18: K=(8x-19), M=(5x+43)}
\text{opposite in quad: sum=180}
\text{8x-19+5x+43=180}
\text{13x = 156, x = 12}
\text{M = 5(12)+43 = 103}


\text{=================================}

\text{TOPIC 7: TANGENTS}

\text{Tangent PERPENDICULAR to radius}
\text{at point of tangency}

\text{CHECK IF TANGENT: Pythagorean}
\text{r^2 + tangent^2 = hypotenuse^2}

\text{If equation holds, it's tangent}

\text{Example Q19: r=4.8, seg=7.2,}
\text{hyp CB=12}
\text{4.8^2 + 7.2^2 = 23.04+51.84=74.88}
\text{12^2 = 144}
\text{NOT equal \to NOT tangent}

\text{Example Q20: r=6.8, seg=15,}
\text{hyp=11.2 (with 15)}
\text{check if 6.8^2+15^2 = ?^2}

\text{TWO TANGENTS from same point}
\text{are CONGRUENT}

\text{Q21: tangent segments equal}
\text{Use Pythagorean on right triangle}
\text{8^2 + 11^2 = (8+x)^2 typically}
\text{Check figure for setup}

\text{Q23: DE = EF (both tangent from E)}
\text{10x-19 = 4x+23}
\text{6x = 42, x = 7}
\text{EF = 4(7)+23 = 51}

\text{Q24: PERIMETER of triangle}
\text{with inscribed circle}
\text{Tangent segments from each vertex}
\text{are equal. Add all sides.}


\text{=================================}

\text{TOPIC 8: ANGLES FROM CHORDS,}
\text{SECANTS, TANGENTS}

\text{INSIDE circle (two chords cross):}

\text{angle = (arc1 + arc2) / 2}

\text{(average of the two arcs)}

\text{OUTSIDE circle (2 secants, 2}
\text{tangents, or secant+tangent):}

\text{angle = (far arc - near arc) / 2}

\text{ON circle (tangent-chord):}

\text{angle = arc / 2}

\text{Example Q25: arcs 92 and 140}
\text{chords cross inside}
\text{AED = (92+140)/2 = 116}
\text{DEC = 180 - 116 = 64}

\text{Example Q27: arcs 51 and 125}
\text{two secants from outside}
\text{KLM = (125-51)/2 = 37}

\text{Example Q28: arcs 136 and 68}
\text{XYZ = (136-68)/2 = 34}

\text{Example Q30: arc 35 and 88}
\text{chords cross: angle at V}
\text{if 88 is the angle, then}
\text{88 = (35 + RU)/2}
\text{176 = 35 + RU}
\text{RU = 141}


\text{=================================}

\text{TOPIC 9: SEGMENT LENGTHS}

\text{TWO CHORDS intersect inside:}

\text{a \cdot b = c \cdot d}

\text{(parts of each chord multiply}
\text{to equal products)}

\text{Example Q37: 30 \cdot 18 = 27 \cdot x}
\text{540 = 27x}
\text{x = 20}

\text{TWO SECANTS from outside:}

\text{whole \cdot outside = whole \cdot outside}

\text{(outside seg)(entire seg) both sides}

\text{Example Q39: 7 \cdot (7+15) = 6 \cdot (6+x)}
\text{(numbers vary - use figure)}

\text{SECANT-TANGENT:}

\text{tangent^2 = whole \cdot outside}

\text{Example Q38: tangent 14, secant}
\text{14^2 = 16 \cdot (16 + x+11)... set up}

\text{Example Q40: tangent=x, sec=12+18}
\text{x^2 = 12(12+18) = 12(30) = 360}
\text{x = \sqrt{360} = 6\sqrt{10} \approx 18.97}


\text{=================================}

\text{TOPIC 10: EQUATIONS OF CIRCLES}

\text{STANDARD FORM:}

\text{(x-h)^2 + (y-k)^2 = r^2}

\text{Center: (h,k)}

\text{Radius: r = \sqrt{right side}}

\text{WATCH SIGNS: flip h and k}

\text{Example Q43: (x+2)^2+(y-7)^2=16}
\text{Center: (-2, 7)}
\text{Radius: \sqrt{16} = 4}

\text{Example Q44: x^2+(y+6)^2=121}
\text{Center: (0,-6)}
\text{Radius: 11, Diameter: 22}

\text{WRITE EQUATION from center + r:}

\text{Q47: center (-3,4), r=7}
\text{(x+3)^2 + (y-4)^2 = 49}

\text{Q48: center (-9,0), d=20, r=10}
\text{(x+9)^2 + y^2 = 100}

\text{Q50: center (12,5), r=\sqrt{89}}
\text{(x-12)^2 + (y-5)^2 = 89}

\text{FROM CIRCUMFERENCE:}

\text{Q51: center (2,-2), C=12\pi}
\text{r=6, r^2=36}
\text{(x-2)^2 + (y+2)^2 = 36}

\text{FROM AREA:}

\text{Q52: center (0,10), A=225\pi}
\text{r^2=225, r=15}
\text{x^2 + (y-10)^2 = 225}

\text{FROM POINT ON CIRCLE:}

\text{r = distance from center to point}

\text{r = \sqrt{(x2-x1)^2+(y2-y1)^2}}

\text{Q53: center (-4,7), point (-1,9)}
\text{r^2 = (-1+4)^2+(9-7)^2 = 9+4=13}
\text{(x+4)^2 + (y-7)^2 = 13}

\text{FROM ENDPOINTS OF DIAMETER:}

\text{Center = midpoint}

\text{midpoint = ((x1+x2)/2, (y1+y2)/2)}

\text{radius = half the distance}

\text{Q54: (-18,-3) and (8,11)}
\text{center = (-5, 4)}
\text{r^2 = (8+5)^2 + (11-4)^2}
\text{= 169 + 49 = 218}
\text{(x+5)^2 + (y-4)^2 = 218}


\text{=================================}

\text{COMPLETING THE SQUARE}

\text{Convert general to standard form}

\text{For x^2 + y^2 + Dx + Ey + F = 0:}

\text{Group x terms, y terms}

\text{Add (D/2)^2 to both sides}

\text{Add (E/2)^2 to both sides}

\text{Factor into (x-h)^2 + (y-k)^2}

\text{Example Q55: x^2+y^2-8x+12y+34=0}

\text{Move 34: x^2-8x+y^2+12y = -34}

\text{(-8/2)^2 = 16}
\text{(12/2)^2 = 36}

\text{x^2-8x+16 + y^2+12y+36 = -34+16+36}

\text{(x-4)^2 + (y+6)^2 = 18}

\text{Center: (4,-6), r=\sqrt{18}=3\sqrt{2}}

\text{Example Q56: x^2+y^2+x=279+3x-6y}

\text{Rearrange: x^2-2x + y^2+6y = 279}

\text{(-2/2)^2 = 1}
\text{(6/2)^2 = 9}

\text{x^2-2x+1 + y^2+6y+9 = 279+1+9}

\text{(x-1)^2 + (y+3)^2 = 289}

\text{Center: (1,-3), r=17}


\text{=================================}

\text{QUICK REFERENCE CHART}

\text{ANGLE LOCATION \to FORMULA}

\text{At center: angle = arc}

\text{On circle (inscribed):}
\text{angle = arc/2}

\text{On circle (tangent-chord):}
\text{angle = arc/2}

\text{Inside (chords cross):}
\text{angle = (sum of arcs)/2}

\text{Outside (2 secants/tangents):}
\text{angle = (diff of arcs)/2}


\text{SEGMENT PRODUCTS}

\text{Chord-chord: a \cdot b = c \cdot d}

\text{Secant-secant:}
\text{whole \cdot outside = whole \cdot outside}

\text{Tangent-secant:}
\text{tan^2 = whole \cdot outside}


\text{TRANSFORMATION RULES}

\text{Translate: (x+a, y+b)}

\text{Reflect x: (x,-y)}
\text{Reflect y: (-x,y)}
\text{Reflect y=x: (y,x)}

\text{Rotate 90 CCW: (-y,x)}
\text{Rotate 180: (-x,-y)}
\text{Rotate 270 CCW: (y,-x)}

\text{Dilate at origin: (kx,ky)}

\text{Dilate at (a,b):}
\text{(k(x-a)+a, k(y-b)+b)}
